\section{Post-result target diagnostics}\label{app:diagnostics}
These retrospective analyses retain the completed 1,008 generated images and
the original inferential family. The initial diagnostics use unchanged reference
vectors and scaling; the subsequent source audit separately varies painting
regions, development scaling and content labels.

\subsection{Content-specific targets and genuine-painting controls}
The recorded reference classes are water-organized, built-place-organized,
route-organized, and open or wooded land, assigned by a fixed source-title
lexicon. Counts in that order are $(191,67,8,31)$ for Monet, $(52,29,9,16)$ for
Sisley, $(28,77,12,24)$ for Pissarro and $(31,28,12,34)$ for C\'ezanne.
Class-specific means are centered across painters just as in the pooled target.
Their mean squared departure from the pooled target is $.438H$, a description
of observed reference means that includes finite-sample error.

For the generated comparison, the four water, four built and three land briefs
use the corresponding class means; the three mixed briefs are excluded.
We average the class-specific cross-product error over these $S=11$ scenes and divide by
$H_c=S^{-1}\sum_s\sum_a\|r_{a,c(s)}\|^2=7.413$, compared with pooled $H=5.915$.
Each column of Table~\ref{tab:review-targets} therefore has its own denominator.
The unchanged 11-scene pooled comparison isolates the scene-subset choice.
These title-derived strata have 16--191 reference works per artist/class among
the three included classes, but do not supply visually matched historical scenes.

For each of 1,000 genuine-painting controls, works in each artist/class cell
are randomly divided into disjoint halves, with the smaller half used for
target estimation. A pooled control samples two independent works with replacement
from each painter's held-out half for each of 14 pseudo-scenes. A class-stratified
control samples within the held-out class for each of the 11 water, built or
land pseudo-scenes. It is compared with both pooled and class-specific means
from the training halves, using each target's own squared magnitude.
The random seed is 2026091201. Sampling can repeat a held-out work, but no
held-out work is used to estimate that split's target. The development scaler
is fixed throughout.

To separate sparse sampling from changes in the comparison pools, we also
calculate each split's expected score from its exact held-out means, keeping
the training target and normalizer fixed (Table~\ref{tab:painting-control-ranges}).

\begin{table}[htbp]
\centering\small
\begin{tabular}{llcc}
\toprule
Control sampling & Reference target & Two draws & Exact held-out means\\
\midrule
Pooled & Pooled & $[-.758,1.380]$ & $[.135,.428]$\\
Class-stratified & Pooled & $[-.554,2.400]$ & $[.544,.993]$\\
Class-stratified & Class-specific & $[-.248,1.946]$ & $[.517,.954]$\\
\bottomrule
\end{tabular}
\caption{Central 95\% ranges of genuine-painting control scores across 1,000
splits. Exact means integrate out sampling within each split; these finite-pool
ranges are not confidence intervals.}
\label{tab:painting-control-ranges}
\end{table}

For pooled controls, SD falls from $.561$ to $.076$ after integrating out
sampling. These controls still condition on imperfect title classes and a
single source capture per work; they do not calibrate artistic fidelity.

\begin{table}[htbp]\centering\small
\begin{tabular}{@{}lrrrr@{}}\toprule
 & \multicolumn{2}{c}{11 common briefs} & \multicolumn{2}{c}{14 briefs}\\
Model & Pooled target & Class target & Equal family & Covariance\\\midrule
GPT Image 1 & 1.557 & 1.366 & 1.551 & 2.103\\
GPT Image 2 & 1.099 & 1.130 & 1.243 & 1.511\\
2.5 Flare & 1.726 & 1.539 & 1.765 & 1.636\\
2.5 Sunburst & 1.633 & 1.525 & 1.680 & 1.617\\
Nano Banana 2 & 0.987 & 0.948 & 1.203 & 1.385\\
FLUX.2 Max & 0.750 & 0.786 & 0.808 & 0.928\\
\bottomrule\end{tabular}
\caption{Descriptive error under alternative targets and metrics. The class target uses title-derived water, built and land strata; the three mixed briefs are excluded. Each column uses its own reference normalizer; absolute values across columns are not directly comparable. Covariance is fitted only on development works, with fixed 50\% shrinkage.}
\label{tab:review-targets}\end{table}

\FloatBarrier

\subsection{Magnitude, artist coverage and coordinate weighting}\label{app:magnitude-coverage}
The cross-repeat version of Equation~\ref{eq:scene-decomposition} is exact for
the retained vectors. Its variation term replaces squared deviations with
cross-products between scene-centered repeats. The term can be negative in
finite samples, although its independent-repeat population target is nonnegative.
Scalar calibration fits only the other 13 scenes. All fold denominators $Q$
are positive. GPT Image 2's fitted scalars range from .438 to .459 and FLUX's
from .598 to .694; all fold-specific scores are retained. No calibration
parameter is selected by minimizing its own held-out score. The subsequent
deletion check refits the entire procedure on
each remaining 13-scene panel, rather than treating overlapping folds as
independent replications. Its mean GPT Image 2-minus-FLUX error differences
range from $-.189$ to $-.137$.

\begin{table}[htbp]\centering\small
\begin{tabular}{@{}lrr@{}}\toprule
Model & Squared contrast magnitude $Q$ & Alignment ratio $\beta/\sqrt Q$\\
\midrule
GPT Image 1 & 2.449 & 0.600\\
GPT Image 2 & 2.223 & 0.670\\
2.5 Flare & 2.281 & 0.511\\
2.5 Sunburst & 2.289 & 0.545\\
Nano Banana 2 & 0.994 & 0.444\\
FLUX.2 Max & 0.741 & 0.546\\
\bottomrule\end{tabular}
\caption{Descriptive magnitude and alignment of painter contrasts. $Q=1$ equals the reference squared size, and $D=1-2\beta+Q$. The alignment ratio can fall outside $[-1,1]$ in finite samples.}
\label{tab:review-alignment}\end{table}


For artist pairs, both generated and reference contrasts are direct differences
between the two artists; the denominator is the squared reference-pair distance.
Leaving one artist out recenters the remaining three generated and reference
conditions. Reference singular components are the three rank-one terms of the
centered four-by-31 reference matrix: each combines an artist contrast with a
feature direction, rather than representing one artist. Component aligned
responses project onto each term separately.

\begin{table}[htbp]\centering\small
\begin{tabular}{@{}lrrr|rrrr@{}}\toprule
 & \multicolumn{3}{c}{Component aligned response} & \multicolumn{4}{c}{Aligned response after omitting}\\
Model & First & Second & Third & Monet & Sisley & Pissarro & C\'ezanne\\\midrule
GPT Image 1 & 1.248 & 0.310 & 0.356 & 1.140 & 1.072 & 0.840 & 0.389\\
GPT Image 2 & 1.156 & 0.617 & 0.800 & 1.195 & 0.972 & 0.992 & 0.651\\
2.5 Flare & 1.022 & 0.211 & 0.383 & 1.034 & 0.734 & 0.802 & 0.222\\
2.5 Sunburst & 1.082 & 0.250 & 0.417 & 1.116 & 0.840 & 0.762 & 0.284\\
Nano Banana 2 & 0.513 & 0.371 & 0.195 & 0.569 & 0.446 & 0.388 & 0.276\\
FLUX.2 Max & 0.666 & 0.021 & 0.183 & 0.539 & 0.511 & 0.527 & 0.096\\
\bottomrule\end{tabular}
\caption{Descriptive artist coverage of the aligned response. Reference components account for 66.3\%, 20.6\% and 13.0\% of centroid variation. Each omission recomputes the target and recenters the other three artists.}
\label{tab:review-components}\end{table}


Table~\ref{tab:review-components} shows that aggregate positive alignment is
uneven across those components.

Equal-family weighting assigns each coordinate inverse family size (11, 8 or 12).
The covariance sensitivity uses 221 development works with equal painter mass.
For their weighted covariance $\Sigma$, the metric is
$[.5\Sigma+.5\operatorname{tr}(\Sigma)I/31]^{-1}$, with fixed shrinkage and no
tuning on evaluated images. These weighting choices are not independently
validated artistic metrics.

\begin{samepage}
FLUX retains the lowest point estimate under all
31 single-coordinate deletions; the report retains every coordinate contribution,
including negative terms.
Without texture, Flare and Sunburst still have error estimates above the
no-contrast value of one (1.580 and 1.604).
\end{samepage}

\subsection{A design-aware energy correction}
The primary energy values use empirical distributions, including their zero
diagonal distances. To estimate energy between a fixed empirical reference
and the equally weighted mixture of scene-conditional output distributions,
cross-scene generated pairs already have the required independent draws.
Only the same-scene generated term needs adjustment. With two repeats and
$S$ scenes, the descriptive correction is
\[
 \widehat{\mathcal E}_{\rm strata}
 =\widehat{\mathcal E}_{\rm empirical}
 -\frac{1}{2S^2}\sum_s\|y_{s1}-y_{s2}\|.
\]
The reference within-distance remains empirical because this target conditions
on the finite reference panel. A pooled IID $n/(n-1)$ multiplier would also
alter cross-scene terms and estimate a different target. This correction retains
the 21 negative named-minus-generic differences. It does not solve unequal
historical and generated content mixtures, and no new energy tests are performed.

\subsection{Uncertainty and synthetic coverage checks}
The Student procedure uses 14 scene summaries. For independent scene errors
with fixed conditional means $\delta_s$, its estimated variance has expectation
\[
 \mathbb E[s_\delta^2/S]
 =\operatorname{Var}(\bar{\widehat\delta})
 +\frac{1}{S(S-1)}\sum_s(\delta_s-\bar\delta)^2.
\]
It therefore includes real between-scene heterogeneity as well as request noise.
This explains why it is not a direct estimate of fresh-request variance for
exactly the fixed panel. It also does not provide an exact coverage guarantee:
non-Gaussian cross-products and dependence can affect the Student approximation.

The initial 5,000-run scenarios (seed 2026091202) use the six configurations,
14 scenes, two repeats and 21-endpoint family, with unit reference magnitude.
Gaussian, variance-standardized $t_3$/heteroskedastic and fixed scene-interaction
errors give simultaneous coverage 95.96\%, 97.42\% and 99.68\%. Their independent
centered noise power is $.5$. Allocating half the noise to a model-specific
state shared by all scenes and both repeats reduces coverage to 0.04\%, with
mean corrected-error bias near $.25$.

A subsequent sweep uses independent Gaussian errors at powers $.5$, $1.5$ and
$3$, with either isotropic or rank-three covariance aligned with the reference
components (5,000 runs per case). Coverage ranges from 95.06\% to 96.80\%, with
Monte Carlo SE $.25$--$.31$ percentage points. For comparison, directly centered
repeat differences estimate observed powers of $.424$--$2.596$ relative to $H$.
These stress models probe coverage approximations, not actual service independence;
no significance threshold is selected from them.

\subsection{Source-quality sensitivity}\label{app:reference-quality}
AI assistants inspect all 649 reference and 221 development reproductions
using contact sheets and enlarged peripheral regions, recording removable
calibration targets, frames, margins and labels. Uncertain boundaries remain
unchanged. Reference content is visually classified before comparison with the
title label; mixed or unclear cases retain that label.

The audit selects crops for 90 reference and 41 development images, including
ten calibration-strip cases; 43 uncertain boundaries remain unchanged. Visual
classes differ from 230 title assignments, while 42 mixed or unclear cases
retain their original labels. These AI-coded sensitivities did not undergo human verification and are not
independent ground truth.

Audited rectangular crops follow the original orientation and ICC rules, before
aspect-preserving Lanczos resizing to a 512-pixel short side without upsampling.
All 31 features and all works are retained. We compare corrected references with
original scaling, then refit development IQR scaling and consistently transform
generated vectors. Label sensitivity uses the same 11 briefs. Versioned records
preserve audit decisions, new vectors and original results; model scores do not
select corrections.
All six aligned-response intervals remain above zero with either scaler.
With original scaling, corrected regions give FLUX $D=0.832$.
After refitting, (uncalibrated, held-out calibrated) errors are:
GPT Image 1 $(1.631,0.693)$, GPT Image 2 $(1.259,0.579)$,
Flare $(1.737,0.760)$, Sunburst $(1.577,0.711)$, Nano Banana 2 $(1.121,0.846)$
and FLUX $(0.872,0.762)$. FLUX and GPT Image 2 retain the respective minima.

\section{Image selection and provenance}\label{app:images}
The generated examples use the first declared brief and first repeat, selected
after analysis without appearance- or score-based curation. The brief reads:
``A broad river bends past a gravel bank. Three tall trees stand on the far bank,
with a low hill behind them under a pale overcast sky.'' Every prompt ends with
``No text or frame.'' Only the clauses in Section~\ref{sec:design} change.

References in Figure~\ref{fig:original-generated} are the lexicographically first
measured work IDs per painter classified as water-organized in the completed
audit, with recorded public-domain status. The Sisley and Pissarro reproductions
use audited crops removing calibration strips; Monet and C\'ezanne retain their
full regions. The manifest records image identities, provenance, crop boxes
and exact prompts. The paintings are not matched to the generated brief.


\FloatBarrier
